% GENERATED FILE. Edit canonical lessons, then run npm run generate:books.
% canonical-source-hash: fnv1a64:1129368879b84527
% canonical-lessons: KA-C247-dasara, KA-C247-dipavali, KA-C247-sankranti, KA-C247-ganesa-caturthi, KA-C247-rajyotsava

\chapter{Dasara, Deepavali, Sankranti, Ganesha Chaturthi, Rajyotsava}
\label{ch:ka-dasara-deepavali-sankranti-ganesha-chaturthi-rajyotsava}
\hlchaptermodality{247}

\begin{chapteropening}
\textbf{By the end of this chapter:} \emph{I can say Dasara, Deepavali, Sankranti, Ganesha Chaturthi and Rajyotsava.}

Complete the last lesson of chapter 247: I can say Dasara, Deepavali, Sankranti, Ganesha Chaturthi and Rajyotsava.
\end{chapteropening}

\section[\texorpdfstring{\kn{ದಸರಾ} (dasarā)}{ದಸರಾ (dasarā)}]{\kn{ದಸರಾ} (dasarā) — Dasara, the autumn festival}

\label{lesson:KA-C247-dasara}



\begin{quote}
Before the new one: say the Kannada for a bandage, then the Kannada for diabetes.
\end{quote}

\subsection*{You'll want to know: \kn{ದಸರಾ}}
\textbf{\kn{ದಸರಾ}} — \emph{dasarā} — \textquotedblleft{}Dasara, the autumn festival\textquotedblright{}.

\begin{grammarlens}[title={What you've built}]
One more word for this chapter.
\end{grammarlens}

\subsection*{Your turn}
\begin{itemize}
\raggedright
  \item \emph{Say it:} \emph{dasarā}
  \item \emph{Say it:} \emph{dasarā}, once more
  \item \emph{Say it:} say \emph{dasarā}
  \item \emph{Recall:} say the Kannada for a bandage, then the Kannada for diabetes, then say \emph{dasarā} again
\end{itemize}

\begin{tcolorbox}[breakable,colback=teal!4,colframe=teal!35!black,title={Before you move on}]
What does \kn{ದಸರಾ} mean? (\textquotedblleft{}Dasara, the autumn festival\textquotedblright{}.) Say it once more.
\end{tcolorbox}
\section[\texorpdfstring{\kn{ದೀಪಾವಳಿ} (dīpāvaḷi)}{ದೀಪಾವಳಿ (dīpāvaḷi)}]{\kn{ದೀಪಾವಳಿ} (dīpāvaḷi) — Deepavali, the festival of lights}

\label{lesson:KA-C247-dipavali}



\begin{quote}
Before the new one: say the Kannada for diabetes, then the Kannada for Dasara.
\end{quote}

\subsection*{You'll want to know: \kn{ದೀಪಾವಳಿ}}
\textbf{\kn{ದೀಪಾವಳಿ}} — \emph{dīpāvaḷi} — \textquotedblleft{}Deepavali, the festival of lights\textquotedblright{}.

\begin{grammarlens}[title={What you've built}]
One more word for this chapter.
\end{grammarlens}

\subsection*{Your turn}
\begin{itemize}
\raggedright
  \item \emph{Say it:} \emph{dīpāvaḷi}
  \item \emph{Say it:} \emph{dīpāvaḷi}, once more
  \item \emph{Say it:} say \emph{dīpāvaḷi}
  \item \emph{Recall:} say the Kannada for diabetes, then the Kannada for Dasara, then say \emph{dīpāvaḷi} again
\end{itemize}

\begin{tcolorbox}[breakable,colback=teal!4,colframe=teal!35!black,title={Before you move on}]
What does \kn{ದೀಪಾವಳಿ} mean? (\textquotedblleft{}Deepavali, the festival of lights\textquotedblright{}.) Say it once more.
\end{tcolorbox}
\section[\texorpdfstring{\kn{ಸಂಕ್ರಾಂತಿ} (saṅkrānti)}{ಸಂಕ್ರಾಂತಿ (saṅkrānti)}]{\kn{ಸಂಕ್ರಾಂತಿ} (saṅkrānti) — Sankranti, the January harvest festival}

\label{lesson:KA-C247-sankranti}



\begin{quote}
Before the new one: say the Kannada for Dasara, then the Kannada for Deepavali.
\end{quote}

\subsection*{You'll want to know: \kn{ಸಂಕ್ರಾಂತಿ}}
\textbf{\kn{ಸಂಕ್ರಾಂತಿ}} — \emph{saṅkrānti} — \textquotedblleft{}Sankranti, the January harvest festival\textquotedblright{}.

\begin{grammarlens}[title={What you've built}]
One more word for this chapter.
\end{grammarlens}

\subsection*{Your turn}
\begin{itemize}
\raggedright
  \item \emph{Say it:} \emph{saṅkrānti}
  \item \emph{Say it:} \emph{saṅkrānti}, once more
  \item \emph{Say it:} say \emph{saṅkrānti}
  \item \emph{Recall:} say the Kannada for Dasara, then the Kannada for Deepavali, then say \emph{saṅkrānti} again
\end{itemize}

\begin{tcolorbox}[breakable,colback=teal!4,colframe=teal!35!black,title={Before you move on}]
What does \kn{ಸಂಕ್ರಾಂತಿ} mean? (\textquotedblleft{}Sankranti, the January harvest festival\textquotedblright{}.) Say it once more.
\end{tcolorbox}
\section[\texorpdfstring{\kn{ಗಣೇಶ} \kn{ಚತುರ್ಥಿ} (gaṇēśa caturthi)}{ಗಣೇಶ ಚತುರ್ಥಿ (gaṇēśa caturthi)}]{\kn{ಗಣೇಶ} \kn{ಚತುರ್ಥಿ} (gaṇēśa caturthi) — Ganesha Chaturthi}

\label{lesson:KA-C247-ganesa-caturthi}



\begin{quote}
Before the new one: say the Kannada for Deepavali, then the Kannada for Sankranti.
\end{quote}

\subsection*{You'll want to know: \kn{ಗಣೇಶ} \kn{ಚತುರ್ಥಿ}}
\textbf{\kn{ಗಣೇಶ} \kn{ಚತುರ್ಥಿ}} — \emph{gaṇēśa caturthi} — \textquotedblleft{}Ganesha Chaturthi\textquotedblright{}.

\begin{grammarlens}[title={What you've built}]
One more word for this chapter.
\end{grammarlens}

\subsection*{Your turn}
\begin{itemize}
\raggedright
  \item \emph{Say it:} \emph{gaṇēśa caturthi}
  \item \emph{Say it:} \emph{gaṇēśa caturthi}, once more
  \item \emph{Say it:} say \emph{gaṇēśa caturthi}
  \item \emph{Recall:} say the Kannada for Deepavali, then the Kannada for Sankranti, then say \emph{gaṇēśa caturthi} again
\end{itemize}

\begin{tcolorbox}[breakable,colback=teal!4,colframe=teal!35!black,title={Before you move on}]
What does \kn{ಗಣೇಶ} \kn{ಚತುರ್ಥಿ} mean? (\textquotedblleft{}Ganesha Chaturthi\textquotedblright{}.) Say it once more.
\end{tcolorbox}
\section[\texorpdfstring{\kn{ರಾಜ್ಯೋತ್ಸವ} (rājyōtsava)}{ರಾಜ್ಯೋತ್ಸವ (rājyōtsava)}]{\kn{ರಾಜ್ಯೋತ್ಸವ} (rājyōtsava) — Rajyotsava, Karnataka Formation Day}

\label{lesson:KA-C247-rajyotsava}



\begin{quote}
Before the new one: say the Kannada for Sankranti, then the Kannada for Ganesha Chaturthi.
\end{quote}

\subsection*{You'll want to know: \kn{ರಾಜ್ಯೋತ್ಸವ}}
\textbf{\kn{ರಾಜ್ಯೋತ್ಸವ}} — \emph{rājyōtsava} — \textquotedblleft{}Rajyotsava, Karnataka Formation Day\textquotedblright{}.

\begin{grammarlens}[title={What you've built}]
One more word for this chapter.
\end{grammarlens}

\subsection*{Your turn}
\begin{itemize}
\raggedright
  \item \emph{Say it:} \emph{rājyōtsava}
  \item \emph{Say it:} \emph{rājyōtsava}, once more
  \item \emph{Say it:} say \emph{rājyōtsava}
  \item \emph{Recall:} say the Kannada for Sankranti, then the Kannada for Ganesha Chaturthi, then say \emph{rājyōtsava} again
\end{itemize}

\begin{tcolorbox}[breakable,colback=teal!4,colframe=teal!35!black,title={Before you move on}]
What does \kn{ರಾಜ್ಯೋತ್ಸವ} mean? (\textquotedblleft{}Rajyotsava, Karnataka Formation Day\textquotedblright{}.) Say it once more.
\end{tcolorbox}
